\chapter*{Tóm tắt}
\addcontentsline{toc}{chapter}{Tóm tắt}
Trích xuất thông tin chính KIE (Key Information Extraction) từ các tài liệu giàu thị giác VRD (Visually Rich Documents) như biểu mẫu hành chính, hóa đơn thương mại và biên lai bán lẻ đóng vai trò nền tảng trong lĩnh vực Trí tuệ nhân tạo cho Tài liệu Document AI và Tự động hóa quy trình bằng robot RPA (Robotic Process Automation). Các mô hình Transformer đa phương thức tiền huấn luyện, tiêu biểu là LayoutLMv3, đã đạt được những bước tiến lớn nhờ kết hợp đồng thời văn bản, tọa độ bố cục hai chiều (2D) và đặc trưng hình ảnh thông qua cơ chế che mặt nạ hợp nhất. Tuy nhiên, kiến trúc gán nhãn chuỗi ở mức token tiêu chuẩn của LayoutLMv3 vẫn tồn tại ba hạn chế cốt lõi về mặt cấu trúc và ngữ cảnh: (1) sự phụ thuộc hoàn toàn vào tọa độ hộp giới hạn 2D tuyệt đối khiến mô hình phải tự suy luận ngầm định các cấu trúc phân cấp vĩ mô như dòng văn bản, khối cục bộ và cột dọc; (2) việc phân loại độc lập trên từng token thiếu sự trao đổi ngữ cảnh tường minh giữa các phân đoạn lân cận theo thứ tự đọc, dẫn tới hiện tượng trôi nhãn BIO trên các thực thể dài nhiều token; và (3) sự thiếu vắng ràng buộc ranh giới tường minh giữa các cặp token liên tiếp khiến mô hình dễ phân mảnh hoặc gộp sai các thực thể liền kề trên cùng một dòng.

Nhằm giải quyết đồng thời các thách thức trên mà không cần tiền huấn luyện lại từ đầu, khóa luận nghiên cứu và đề xuất khung kiến trúc cải tiến trên nền tảng LayoutLMv3 (triển khai trong lớp mô hình \texttt{LayoutLMv3For}\allowbreak\texttt{SegmentTokenClassification}) tích hợp ba thành phần bổ trợ lẫn nhau: (1) \textbf{Mã hóa vị trí phân cấp HPE (Hierarchical Position Encoding)} trích xuất chỉ số dòng, khối và cột bằng các thuật toán gom cụm hình học để nhúng trực tiếp cấu trúc bố cục vĩ mô vào biểu diễn đầu vào; (2) \textbf{Mô-đun ngữ cảnh phân đoạn SCM (Segment Context Module)} thực hiện gộp trung bình các token trong cùng phân đoạn, lan truyền ngữ cảnh liên phân đoạn qua bộ mã hóa Transformer nhẹ kết hợp cổng dư ReZero và nhúng đánh dấu token đầu phân đoạn (\texttt{is\_first\_}\allowbreak\texttt{token\_embedding}) nhằm duy trì tính nhất quán ngữ nghĩa đồng thời bảo toàn chuẩn gán nhãn BIO; và (3) \textbf{Hàm mất mát ranh giới nội dòng/thực thể} huấn luyện bộ phân loại nhị phân trên từng cặp token liên tiếp dựa trên định danh thực thể liên tục (\texttt{entity\_ids}) để củng cố ranh giới chuyển tiếp thực thể. Thực nghiệm kiểm chứng nghiêm ngặt qua 3 hạt giống ngẫu nhiên $\{42, 123, 1993\}$ trên máy chủ GPU NVIDIA A100-SXM4-80GB với hai bộ dữ liệu chuẩn quốc tế FUNSD và CORD cho thấy mô hình đề xuất cải thiện vượt trội điểm F1 mức thực thể từ $89.92 \pm 0.85\%$ lên $91.58 \pm 0.16\%$ ($+1.66\%$, giảm hơn 5 lần độ lệch chuẩn) trên tập FUNSD với biến thể \texttt{base}, đồng thời duy trì độ chính xác và tính ổn định cao trên tập dữ liệu biên lai CORD.

\textbf{\textit{Từ khóa:}} Trích xuất thông tin chính KIE (Key Information Extraction), Tài liệu giàu thị giác VRD (Visually Rich Documents), LayoutLMv3, Mã hóa vị trí phân cấp HPE (Hierarchical Position Encoding), Mô-đun ngữ cảnh phân đoạn SCM (Segment Context Module), Hàm mất mát ranh giới, Gán nhãn chuỗi BIO.

